%
% k2.tex -- template for standalon tikz images
%
% (c) 2021 Prof Dr Andreas Müller, OST Ostschweizer Fachhochschule
%
\documentclass[tikz]{standalone}
\usepackage{amsmath}
\usepackage{times}
\usepackage{txfonts}
\usepackage{pgfplots}
\usepackage{csvsimple}
\usetikzlibrary{arrows,intersections,math}
\begin{document}
\def\skala{1}
\begin{tikzpicture}[>=latex,thick,scale=\skala]

\def\r{0.82}
\coordinate (Z)   at (0,0);
\coordinate (M0)  at (0,{\r});
\coordinate (M1)  at (0,{2*\r});
\coordinate (M2)  at (0,{4*\r});
% P1 Punkte (senkrecht zu ZM1 durch M1)
\coordinate (P11) at ({2*\r},{2*\r});
\coordinate (P12) at ({-2*\r},{2*\r});
% P2 Punkte (auf Diagonalen von M2, Radius 4r)
% 2.828 = 4*0.7071 = 4/sqrt2
\coordinate (P21) at ({2.828*\r},{1.172*\r});   % 315 Grad, unten-rechts
\coordinate (P22) at ({2.828*\r},{6.828*\r});   % 45 Grad, oben-rechts
\coordinate (P23) at ({-2.828*\r},{6.828*\r});  % 135 Grad, oben-links
\coordinate (P24) at ({-2.828*\r},{1.172*\r});  % 225 Grad, unten-links
% Kreise
\draw (M0) circle ({\r});
\draw (M1) circle ({2*\r});
\draw (M2) circle ({4*\r});
% Geraden durch M2
\draw[thin] (0,{-0.6}) -- (0,{8*\r+0.6});
\draw[thin] ({-4*\r-0.6},{4*\r}) -- ({4*\r+0.6},{4*\r});
\draw[thin] ({-2.828*\r-0.55},{6.828*\r+0.55}) -- ({2.828*\r+0.55},{1.172*\r-0.55});
\draw[thin] ({2.828*\r+0.55},{6.828*\r+0.55}) -- ({-2.828*\r-0.55},{1.172*\r-0.55});
% ROTE LINIEN: Dreieck 1: P2,3 -- P2,1 -- Z; Dreieck 2: Z -- P1,1 -- P2,1
\draw[red, thick] (P23) -- (P21) -- (Z) -- cycle;
\draw[red, thick] (Z) -- (P11) -- (P21) -- cycle;
% Punkte – Labels mit weissem Hintergrund, kollisionsfrei
\fill (Z)   circle (1.5pt) node[below left=4pt,  fill=white, inner sep=1pt, font=\small] {$Z$};
\fill (M0)  circle (1.5pt) node[right=6pt,       fill=white, inner sep=1pt, font=\small] {$M_0$};
\fill (M1)  circle (1.5pt) node[right=6pt,       fill=white, inner sep=1pt, font=\small] {$M_1$};
\fill (M2)  circle (1.5pt) node[right=6pt,       fill=white, inner sep=1pt, font=\small] {$M_2$};
\fill (P11) circle (1.5pt) node[above right=6pt, fill=white, inner sep=1pt, font=\small] {$P_{1,1}$};
\fill (P12) circle (1.5pt) node[above left=6pt,  fill=white, inner sep=1pt, font=\small] {$P_{1,2}$};
\fill (P21) circle (1.5pt) node[below right=6pt, fill=white, inner sep=1pt, font=\small] {$P_{2,1}$};
\fill (P22) circle (1.5pt) node[above right=6pt, fill=white, inner sep=1pt, font=\small] {$P_{2,2}$};
\fill (P23) circle (1.5pt) node[above left=6pt,  fill=white, inner sep=1pt, font=\small] {$P_{2,3}$};
\fill (P24) circle (1.5pt) node[below left=6pt,  fill=white, inner sep=1pt, font=\small] {$P_{2,4}$};
% k-Labels mit weissem Hintergrund
\node[font=\small, fill=white, inner sep=1pt] at ({-4*\r-0.5},{6.5*\r}) {$k_2$};
\node[font=\small, fill=white, inner sep=1pt] at ({-2.2*\r},{2.0*\r}) {$k_1$};

\end{tikzpicture}
\end{document}

